\chapter{Disease descriptions}
\label{sec:pathologies}

\section{Chest Pathologies in CXR}

Chest X-ray images encode a broad spectrum of pathological findings, ranging from subtle parenchymal opacities 
to gross structural abnormalities. A thorough understanding of these conditions is essential for appreciating 
the challenges posed to both radiologists and automated diagnostic systems. In this work, we focus on 
pathologies annotated in the CheXpert dataset, which we group below by anatomical compartment.

\subsection{Airspace and Parenchymal Diseases}

Airspace diseases are characterized by abnormal filling or collapse of the alveolar spaces, and are among 
the most commonly encountered findings on CXR.

\subsubsection{Atelectasis}
Atelectasis refers to the partial or complete collapse of lung tissue, resulting in reduced or absent aeration 
of the affected region. On CXR, it typically manifests as increased opacity with associated volume loss, 
displacement of fissures, and shift of the mediastinum toward the affected side. It may arise from 
obstruction of a bronchus, external compression, or loss of surfactant.

\subsubsection{Consolidation}
Consolidation occurs when the normally air-filled alveoli are replaced by fluid, cells, or other material, 
resulting in a homogeneous increase in lung density. It is radiologically characterised by airspace opacification 
with air bronchograms, and is commonly associated with pneumonia, pulmonary haemorrhage, or organising 
pneumonia.

\subsubsection{Lung Opacity}
Lung opacity is a broad descriptor encompassing any increase in radiodensity within the pulmonary parenchyma, 
including ground-glass opacity and frank consolidation. It may reflect a range of underlying processes 
including infection, oedema, haemorrhage, and interstitial fibrosis. Automated classification systems often 
encounter difficulty in disambiguating opacity subtypes due to their overlapping appearances.

\subsubsection{Lung Lesion}
Lung lesions refer to discrete focal abnormalities within the pulmonary parenchyma, including nodules, 
masses, and cavitary lesions. They carry significant clinical importance as potential indicators of 
malignancy, tuberculosis, or fungal infection. Due to their relatively small size and low prevalence 
in population-level datasets, they present a significant challenge for deep learning models.

\subsubsection{Pneumonia}
Pneumonia is an infectious process resulting in inflammatory consolidation of one or more lung lobes. 
Radiographic findings include lobar or segmental opacification, air bronchograms, and in advanced cases, 
pleural effusion. Despite being a clinically common diagnosis, it exhibits significant visual overlap with 
other airspace pathologies, contributing to label uncertainty in automated annotation pipelines.

\subsubsection{Edema}
Pulmonary oedema arises from accumulation of fluid in the alveolar and interstitial spaces, most commonly 
secondary to cardiac failure or acute lung injury. CXR findings include bilateral perihilar opacification 
in a ``butterfly'' or ``bat-wing'' distribution, Kerley B lines, and upper lobe vascular redistribution. 
It is frequently co-occurring with cardiomegaly and pleural effusions, introducing multi-label complexity 
in classification tasks.

\subsection{Pleural Diseases}

Pleural pathologies affect the lining of the lungs and chest wall, and are readily identifiable on CXR 
due to the distinct density interfaces they produce.

\subsubsection{Pleural Effusion}
Pleural effusion refers to abnormal accumulation of fluid within the pleural space. On an upright CXR, 
small effusions present as blunting of the costophrenic angle, while larger effusions cause homogeneous 
basal opacification with a meniscal upper border. Pleural effusion is one of the most prevalent and 
reliably labelled findings in large-scale CXR datasets.

\subsubsection{Pleural Other}
This label encompasses less common pleural abnormalities not captured under effusion, including pleural 
thickening, pleural plaques (often associated with asbestos exposure), and fibrothorax. These findings 
are less common in general populations, resulting in significant class imbalance in CXR datasets.

\subsubsection{Pneumothorax}
Pneumothorax is defined as the presence of air in the pleural space, leading to lung collapse. It presents 
as a hyperlucent region at the lung apex devoid of vascular markings, with a visible pleural line. 
Tension pneumothorax, a life-threatening variant, may cause mediastinal shift away from the affected side. 
Automated detection of pneumothorax is clinically critical due to its potential for rapid deterioration.

\subsection{Cardiac and Mediastinal Diseases}

\subsubsection{Cardiomegaly}
Cardiomegaly is defined as an abnormal enlargement of the cardiac silhouette, commonly assessed via 
the cardiothoracic ratio on CXR. A ratio exceeding 0.5 on a posteroanterior (PA) view is generally 
considered diagnostic. It is associated with a range of underlying conditions including congestive heart 
failure, valvular disease, and cardiomyopathy. Cardiomegaly is notable in this work as it is one of 
the pathologies for which latent space visualization was performed in \Cref{chap:VAE}.

\subsubsection{Enlarged Cardiomediastinum}
Enlargement of the cardiomediastinal silhouette encompasses widening of the mediastinum beyond normal 
anatomical limits. This may reflect mediastinal masses, lymphadenopathy, aortic aneurysm, or superior 
vena cava dilation, in addition to cardiomegaly. Distinguishing cardiomegaly from a widened mediastinum 
is often context-dependent and contributes to label ambiguity in automated annotation.

\subsection{Structural and Traumatic Findings}

\subsubsection{Fracture}
Rib and clavicular fractures are identifiable on CXR as cortical discontinuities with or without 
displacement. They are clinically relevant as they may indicate underlying pulmonary contusion, 
pneumothorax, or haemothorax. Fractures represent one of the most infrequent labels in standard 
CXR datasets, contributing to severe class imbalance.

\subsection{Support Devices}

\subsubsection{Support Devices}
This label encompasses the presence of medical devices visible on CXR, including endotracheal tubes, 
central venous catheters, chest drains, pacemaker leads, and nasogastric tubes. Correct identification 
of device placement is clinically important for patient safety. From a machine learning perspective, 
the presence of support devices may confound pathology classification, as critically ill patients 
with devices often concurrently exhibit multiple pathologies.

\subsection{No Finding}

The \textit{No Finding} class denotes CXRs with no identifiable radiographic abnormality. It constitutes 
the largest class in most CXR datasets due to the natural prevalence of healthy subjects in screening 
populations. In automated classification, the dominance of this class relative to rare pathologies 
contributes to significant label imbalance, necessitating careful handling during model training 
and evaluation.

\subsection{Summary and Challenges}

\Cref{tab:pathology_overview} summarises the pathologies described above. A consistent challenge 
across all categories is the co-occurrence of multiple findings within a single image, the subtlety 
of early-stage disease, and the substantial variability in radiographic appearance arising from 
differences in patient positioning, exposure settings, and anatomy. These factors collectively 
motivate the development of robust generative models capable of producing high-fidelity, 
pathology-controlled synthetic CXR images.

\begin{table}[h!]
\centering
\small
\begin{tabular}{llp{7.5cm}}
\toprule
\textbf{Category} & \textbf{Pathology} & \textbf{Key Radiographic Feature} \\
\midrule
\multirow{6}{*}{Airspace / Parenchymal}
    & Atelectasis            & Increased opacity with volume loss \\
    & Consolidation          & Homogeneous opacification, air bronchograms \\
    & Lung Opacity           & Generalised increase in radiodensity \\
    & Lung Lesion            & Discrete focal nodule or mass \\
    & Pneumonia              & Lobar/segmental consolidation \\
    & Edema                  & Bilateral perihilar opacification \\
\midrule
\multirow{3}{*}{Pleural}
    & Pleural Effusion       & Costophrenic blunting, basal opacity \\
    & Pleural Other          & Thickening, plaques, fibrothorax \\
    & Pneumothorax           & Hyperlucent zone, visible pleural line \\
\midrule
\multirow{2}{*}{Cardiac / Mediastinal}
    & Cardiomegaly           & Cardiothoracic ratio $>$ 0.5 \\
    & Enlarged Cardiomediastinum & Widened mediastinal silhouette \\
\midrule
Structural / Traumatic & Fracture & Cortical discontinuity of ribs/clavicle \\
\midrule
Device-related & Support Devices & Visible catheters, tubes, leads \\
\midrule
Normal & No Finding & Absence of radiographic abnormality \\
\bottomrule
\end{tabular}
\caption{Overview of CXR pathologies considered in this work, grouped by anatomical category.}
\label{tab:pathology_overview}
\end{table}